% !TeX root = ../navier-stokes-manuscript.tex
% ============================================================
% 1.2 Where the Problem Is Today
% ============================================================
\subsection{Where the Problem Is Today}\label{sub:HistoricalEstimates}

Local theory gives the smooth solution a beginning.
The first estimate that reaches beyond that local interval is the energy identity.

\newpage
\divider{Energy and the First Uncontrolled Derivative}

For a smooth solution, multiply the momentum equation by \(u\) and integrate over space.
Incompressibility gives
\[
  \int_{\mathbb R^3} (u\cdot\nabla)u\cdot u\,dx
  =
  \frac12\int_{\mathbb R^3} u\cdot\nabla |u|^2\,dx
  =
  0
\]
and
\[
  \int_{\mathbb R^3} \nabla p\cdot u\,dx
  =
  -\int_{\mathbb R^3} p\,\nabla\cdot u\,dx
  =
  0.
\]
Integration by parts turns the viscous term into
\[
  \nu\int_{\mathbb R^3} \Delta u\cdot u\,dx
  =
  -\nu\lVert\nabla u\rVert_{L^2}^{2}.
\]
The three calculations leave
\[
  \lVert u(t)\rVert_{L^2}^{2}
  +
  2\nu\int_0^t\lVert\nabla u(s)\rVert_{L^2}^{2}\,ds
  =
  \lVert u_0\rVert_{L^2}^{2}.
\]
The kinetic energy is bounded at every time, and the squared velocity gradient has a finite integral over every finite time interval.

Leray carried the corresponding inequality through a regularization and compactness argument to construct global weak solutions \parencite{Leray1934}.
On every finite interval they satisfy
\[
  u\in
  L^\infty(0,T;L^2(\mathbb R^3))
  \cap
  L^2(0,T;H^1(\mathbb R^3)).
\]
This gives global existence in the energy class.
Smoothness and uniqueness for arbitrary three-dimensional data remain outside that estimate.

The distinction between the two time controls is already enough to see the opening in Leray's result.
The \(L^2\) norm of \(u\) has a ceiling at every instant.
The \(L^2\) norm of \(\nabla u\) is controlled after integration in time, so high values over short intervals can still have a finite total integral.
The energy identity leaves room for derivative concentration near a finite time.

Vorticity makes the three-dimensional source of that concentration visible.
For
\[
  \omega:=\nabla\times u,
\]
the identity
\[
  \lVert\nabla u\rVert_{L^2}^{2}
  =
  \lVert\omega\rVert_{L^2}^{2}
\]
holds for divergence-free fields with the rapid decay assumed here.
The squared \(L^2\) norm of \(\omega\) is the enstrophy, and its evolution follows from
\[
  \partial_t\omega+(u\cdot\nabla)\omega
  =
  (\omega\cdot\nabla)u+\nu\Delta\omega.
\]
Taking the \(L^2\) inner product with \(\omega\) gives
\[
  \frac12\frac{d}{dt}\lVert\omega\rVert_{L^2}^{2}
  +
  \nu\lVert\nabla\omega\rVert_{L^2}^{2}
  =
  \int_{\mathbb R^3}(\omega\cdot\nabla)u\cdot\omega\,dx.
\]
The integral on the right is the vortex-stretching contribution.
It has no fixed sign, and Leray's estimate gives no pointwise-in-time inequality forcing the diffusive term to dominate it.
Any direct regularity argument has to reach beyond the energy level and control derivatives strongly enough to continue the classical solution.

\divider{What Is Strong Enough to Continue the Solution}

A Sobolev norm measures several spatial derivative orders at once.
For an integer \(m\geq0\),
\[
  \lVert u\rVert_{H^m}^{2}
  :=
  \sum_{|\alpha|\leq m}
  \lVert\partial^\alpha u\rVert_{L^2}^{2},
\]
with the usual Fourier definition for fractional \(m\).
In three dimensions, Sobolev embedding gives \parencite[Theorem~7.28]{Cerda2010}
\[
  H^s(\mathbb R^3)\hookrightarrow C^k(\mathbb R^3)
  \qquad
  \text{whenever}
  \qquad
  s>k+\frac32.
\]
The embedding is the step that turns integral control of spatial derivatives into ordinary bounded continuous derivatives.
In particular,
\[
  s>\frac52
  \qquad\Longrightarrow\qquad
  H^s(\mathbb R^3)\hookrightarrow W^{1,\infty}(\mathbb R^3),
\]
so an \(H^s\) bound above that threshold controls both \(u\) and \(\nabla u\) pointwise.

Suppose a smooth solution on \([0,T)\) satisfies
\[
  \sup_{0\leq t<T}\lVert u(t)\rVert_{H^s}
  \leq M
  \qquad
  \text{for some }s>\frac52.
\]
For every fixed finite \(m\geq s\), differentiating the equation through order \(m\) gives the standard estimate
\[
  \frac12\frac{d}{dt}\lVert u\rVert_{H^m}^{2}
  +
  \nu\lVert\nabla u\rVert_{H^m}^{2}
  \leq
  C_m\lVert\nabla u\rVert_{L^\infty}
  \lVert u\rVert_{H^m}^{2}.
\]
The smooth datum starts with a finite \(H^m\) norm for every finite \(m\).
The assumed \(H^s\) bound controls \(\lVert\nabla u\rVert_{L^\infty}\), and Gronwall's inequality then keeps each chosen \(H^m\) norm finite up to \(T\).
The constant is allowed to depend on \(m\).
Applying Sobolev embedding at each higher finite order recovers every continuous spatial derivative of \(u\).

Pressure comes from the same velocity through
\[
  -\Delta p
  =
  \sum_{i,j=1}^{3}
  \partial_i\partial_j(u_i u_j),
\]
with the usual decay normalization on \(\mathbb R^3\).
Once the spatial derivatives of \(u\) and \(p\) are controlled, the momentum equation gives
\[
  \partial_tu
  =
  \nu\Delta u-(u\cdot\nabla)u-\nabla p.
\]
Repeated differentiation of these two equations recovers every finite mixed spatial and temporal derivative of the velocity--pressure pair.
This is how one sufficiently strong Sobolev estimate reaches the \(C^\infty\) smoothness required by the problem.

The same bound also crosses the alleged endpoint.
Local well-posedness in \(H^s\) gives a positive lifespan determined by \(s\), \(\nu\), and the size of the initial \(H^s\) norm \parencite{FujitaKato1964,Kato1984}.
Choose a time \(t_0<T\) close to \(T\) and use \(u(t_0)\) as new initial data.
The uniform bound by \(M\) gives a common lower bound for the restarted lifespan, so a sufficiently late \(t_0\) carries the restarted solution past \(T\).
Uniqueness identifies it with the original solution wherever their intervals overlap.
Consequently, a finite maximal time \(T_*\) forces
\[
  \lim_{t\uparrow T_*}
  \lVert u(t)\rVert_{H^s}
  =
  \infty
  \qquad
  \text{for every }s>\frac52.
\]
The energy identity supplies \(L^\infty_tL^2_x\) and time-integrated \(H^1_x\) control.
Continuation requires a pointwise-in-time bound at a much higher derivative level.

\divider{Critical Scale and Singular Concentration}

The Navier--Stokes scaling
\[
  u_\lambda(x,t)
  :=
  \lambda u(\lambda x,\lambda^2t),
  \qquad
  p_\lambda(x,t)
  :=
  \lambda^2p(\lambda x,\lambda^2t)
\]
preserves the form of the equations.
A norm whose size is unchanged by this transformation measures concentration at the same strength at every scale.
The Ladyzhenskaya--Prodi--Serrin criteria show how spatial intensity and temporal persistence combine at that critical balance \parencite{Prodi1959,Serrin1962,Serrin1963,Ladyzhenskaya1967}.
They give regularity through \(T\) when
\[
  u\in L^p(0,T;L^q(\mathbb R^3)),
  \qquad
  \frac{2}{p}+\frac{3}{q}\leq1,
  \qquad
  q>3.
\]
The exponent \(q\) measures how concentrated the velocity may become in space, and \(p\) measures how long that concentration may persist.
Equality in the exponent relation is invariant under the scaling above.

At the spatial endpoint,
\[
  \lVert u_\lambda(t)\rVert_{L^3}
  =
  \lVert u(\lambda^2t)\rVert_{L^3}.
\]
Escauriaza, Seregin, and \v{S}ver{\'a}k proved on \(\mathbb R^3\) that
\[
  u\in L^\infty(0,T;L^3(\mathbb R^3))
\]
still gives regularity through \(T\) \parencite{EscauriazaSereginSverak2003}.
Their argument assumes a first singular time, rescales around a concentration, extracts an ancient limiting solution, and uses backward uniqueness to force the limit to vanish.
The endpoint theorem rules out every finite singularity whose critical \(L^3\) norm remains bounded.

Partial regularity gives a local version of the same concentration picture.
For a spacetime point \(z_0=(x_0,t_0)\), let
\[
  Q_r(z_0)
  :=
  B_r(x_0)\times(t_0-r^2,t_0).
\]
Caffarelli, Kohn, and Nirenberg proved an epsilon-regularity statement for suitable weak solutions: sufficiently small scale-invariant velocity and pressure quantities on \(Q_r(z_0)\) force regularity on a smaller cylinder \parencite{CaffarelliKohnNirenberg1982}.
At a singular point, every sufficiently small cylinder must retain a definite amount of scale-invariant concentration.
Their covering argument confines the possible singular set to zero one-dimensional parabolic Hausdorff measure.
Even one point in that set would end global smoothness.

The endpoint \(L^3\) result was later sharpened into a full necessary divergence statement.
For the smooth compactly supported whole-space data in Seregin's theorem, a finite maximal time satisfies \parencite{Seregin2012}
\[
  \lim_{t\uparrow T_*}
  \lVert u(t)\rVert_{L^3(\mathbb R^3)}
  =
  \infty.
\]
The norm eventually exceeds every fixed level and stays above it as \(t\) approaches \(T_*\).
Tao's quantitative version shows that, along infinitely many times approaching a finite endpoint, this norm must grow at least like a positive power of a triple logarithm of the inverse time remaining \parencite{Tao2021Quantitative}.
That lower rate is far slower than any power law, so the known critical estimates still leave a great deal of freedom in how rapidly the concentration can advance.

\divider{What the Estimates Leave Open}

The established results give a concrete description of any finite maximal time.
The kinetic energy stays bounded.
Every continuation-grade Sobolev norm diverges, the whole-space critical \(L^3\) norm diverges, and each singular point retains scale-invariant concentration through every sufficiently small parabolic neighbourhood.

Scaling explains how these statements can coexist.
At corresponding times,
\[
  \lVert u_\lambda(t)\rVert_{L^3}
  =
  \lVert u(\lambda^2t)\rVert_{L^3},
  \qquad
  \lVert u_\lambda(t)\rVert_{L^2}^{2}
  =
  \lambda^{-1}
  \lVert u(\lambda^2t)\rVert_{L^2}^{2}.
\]
Concentration to a shorter length scale preserves the critical \(L^3\) size while reducing the \(L^2\) energy associated with that rescaled structure.
The global energy ceiling gives no scale-independent ceiling on the critical norm.

The parabolic time scale shrinks with the square of the spatial scale.
If a heuristic cascade halves its active length at each stage, the corresponding durations are proportional to
\[
  1,\ \frac14,\ \frac1{16},\ \ldots,
\]
whose sum is finite.
This is the Zeno-cascade picture: known estimates permit a sequence of smaller and faster concentrations to fit below one finite time.
It is a picture of the timing left open by the estimates, and it is not a construction of a Navier--Stokes singularity.

Together these results specify what a finite failure must look like: which quantities must diverge, where concentration must persist, and which continuation bounds would exclude it.
The estimate still missing from the Gold route is one derived from arbitrary smooth data, fixed viscosity, and the equations that keeps a continuation-grade norm finite through every candidate endpoint.
